\histitem{部分算子与谱定理}

一般的希尔伯特空间谱理论在很大程度上是为解决量子力学的数学基础问题而创立的，尤其是在冯·诺伊曼的推动下。

关于量子理论直至 W. 海森堡那篇引入量子力学的奠基性论文 {[}28{]}的发展详述，参见 M. 贾默的工作 {[}35{]}。1924年夏发表该文与1929年发表文章 {[}50{]} 相距不到五年；在后者中，冯·诺伊曼以极为清晰严谨的方式呈现了希尔伯特空间上部分（常称“无界”）埃尔米特算子的全部基本结果。

通过给出希尔伯特空间的第一种公理化表述 {[}49{]}，冯·诺伊曼完成了一个酝酿已久的思想（参见 ÉHM，第267页），并阐明了施密特、弗雷歇、菲舍尔和里斯在1910年前已熟知的序列空间或函数空间之间的各种同构。这种抽象程度的提升使冯·诺伊曼在概念上取得重大进展。

部分算子的思想便是如此；它似乎仍以未成形的形态盘旋在希尔伯特之后重新研究弗雷德霍姆方程的几部著作之上（见上文第1号）。希尔伯特注意到，许多 Sturm--Liouville 型微分方程，特别是在无界区间情形，可化为弗雷德霍姆方程；但是要应用积分方法，必须放宽对核施加的条件，并离开弗雷德霍姆、希尔伯特和施密特的结果适用的框架。多年来，人们研究了越来越多的奇异核，直到开始考虑这样的方程：其左端只有在未知量属于 L²(I) 的子空间时才有意义。在这方面必须提到 E. 希尔布 {[}30{]} 1908年的工作、H. 外尔关于 Sturm--Liouville 理论的1910年工作 {[}82{]}，以及最重要的 T. 卡勒曼 {[}11{]} 1923年的工作，后者在非常一般的框架中展开。

通过统一这些问题，冯·诺伊曼使这些经典工作呈现全新面貌。他尤其强调对称算子与自伴算子之间的根本区别，并用抽象“边界条件”来解释对称算子的自伴扩张，从而推广了在 \(R^{n}\) 的有界开集上拉普拉斯算子的狄利克雷问题和诺伊曼问题中出现的众所周知的边界条件。

与此同时，冯·诺伊曼 {[}51{]} 提出了正规算子 \(^{(5)}\) 的概念，将其作为谱理论的自然框架，并强调由算子及其伴随生成的交换代数的作用。

依托这一精确的数学形式主义——它还揭示了海森堡与薛定谔两种观点的等价性——以及对实验步骤和结果的哥本哈根诠释，非相对论量子力学从此达到一种此后几乎未有变化的精致程度；尽管这套物理理论的结果惊人、甚至看似悖论，其有效性自那以后仍不断得到实验确认，而且精度非凡。

从形式观点看，谱理论的基础由冯·诺伊曼奠定。与此同时，在后者首次提出这些思想之后，M. 斯通于1928年至1930年间进行了类似研究，并获得密切相关的结果。斯通于1932年发表的已知结果的清晰完整表述 {[}72{]}，对希尔伯特谱理论的传播产生了重要影响。

然而，表述上的其他改进也对这些结果的传播很重要。事实上，冯·诺伊曼给出的谱定理（参见 IV，第266页定理1）长期以来被认为非常困难，主要是因为它是在向量值测度框架中表述的。

后来，P. 哈尔莫斯 {[}27{]} 强调通过 \(L^{2}\) 空间上的乘法算子解释谱定理，从而揭示了它在有界算子情形下本质上初等的一面。最近，S. 沃罗诺维奇（在希尔伯特模上的正规算子这一稍有不同的 C*-代数框架中 {[}92{]}）通过引入

“有界化”（参见 IV，第265页第2号），使得可以像处理自伴算子一样简单地处理正规部分算子。

此外，若 D 是 \(R^{n}\) 的开集 U 上的形式对称标量微分算子，则研究 D 在 \(L^{2}(U)\) 上的自伴扩张自然与研究 U 上的广义函数相联系（拉普拉斯算子的情形参见 IV，第242页第6号）。关于后者的历史，参见 J. 迪厄多内的著作 {[}16，第VII章{]}——这是一段曲折的历史，紧接在我们刚才述及的时期之后展开，其中 S. 索伯列夫于1936年 {[}69{]} 和 L. 施瓦茨十年后的决定性贡献值得一提 {[}65{]}。

从纯数学观点看，谱理论与黎曼几何以一种特别富有成果的方式相互作用。因此，M. 卡茨于1966年重新讨论 S. 博赫纳提出的问题，给出了一种引人注目且极其生动的表述 {[}36{]}：确定一个黎曼流形的哪些几何性质可以从拉普拉斯算子的谱性质推出。物理学家 A. 舒斯特于1882年已经预见到这个问题 {[}64{]}：要弄清振动系统发出的不同音调，在某些特殊情形下可能可以解决，也可能不行；但是若要解决逆问题，根据钟能够发出的声音来确定钟的形状，即便是最有技巧的数学家也会束手无策。（“确定振动系统发出的谐波，在某些特殊情形下可能可以解决，也可能不行；但如果要解决逆问题、从钟能够发出的声音确定其形状，最有技巧的数学家也会感到困惑。”）。

